\clearpage
\VerdierPhysicalPageAnchor{0232}%
\begingroup\small\setlength{\parskip}{0pt}\setlength{\emergencystretch}{1.5em}
\setlength{\abovedisplayskip}{4pt plus 2pt minus 2pt}%
\setlength{\belowdisplayskip}{4pt plus 2pt minus 2pt}%

\VerdierSemanticAnchor{ch4d-4-1-7-double-complex-case}
そこで $X$ を二重複体とし、$m=1$ と仮定する。a) を示すには、
$\mathcal{F}\ell_2(\widetilde{\mathbb Z})$ の各 $(p,q,r)$ に対し、
$\operatorname{comp}(\mathcal A)$ の図式

\VerdierSemanticAnchor{ch4d-4-1-7-6}
\begin{equation}
\resizebox{0.98\linewidth}{!}{$
\displaystyle
\int_2{}^{1}\!X(p,q)\longrightarrow
\int_2{}^{1}\!X(p,r)\longrightarrow
\int_2{}^{1}\!X(q,r)
\xrightarrow{{}^{1}\!d(p,q,r)}
\int_2T_1\bigl({}^{1}\!X(p,q)\bigr)
=\left(\int_2{}^{1}\!X(p,q)\right)[1]
$}
\tag{4.1.7.6}
\end{equation}

が $K(\mathcal A)$ において完全三角を定めることを示せば十分である
（この場合、シフトとの可換同型
$\int_2T_1({}^{1}\!X(p,q))\xrightarrow{\sim}
(\int_2{}^{1}\!X(p,q))[1]$ は恒等射であることに注意する）。a) に含まれる
他の主張は明らかである。図式 (4.1.7.6) が $K(\mathcal A)$ において完全
三角を定めることを示すには、性質 (TRII)（第 II 章、1.1.1）により、図式

\VerdierSemanticAnchor{ch4d-4-1-7-7}
\begin{equation}
\resizebox{0.96\linewidth}{!}{$
\displaystyle
\left(\int_2{}^{1}\!X(q,r)\right)[-1]
\xrightarrow{-{}^{1}\!d(p,q,r)[-1]}
\int_2{}^{1}\!X(p,q)\longrightarrow
\int_2{}^{1}\!X(p,r)\longrightarrow
\int_2{}^{1}\!X(q,r)
$}
\tag{4.1.7.7}
\end{equation}

が $K(\mathcal A)$ において完全三角を定めることを示せば十分である。
単純複体 $\int_2{}^{1}\!X(p,r)$ は射 $-{}^{1}\!d(p,q,r)[-1]$ の写像錐
（第 I 章、3.1.2）に等しく、(4.1.7.7) に現れる射
$\int_2{}^{1}\!X(p,q)\longrightarrow\int_2{}^{1}\!X(p,r)$ と
$\int_2{}^{1}\!X(p,r)\longrightarrow\int_2{}^{1}\!X(q,r)$ は、写像錐関手の
考察で導入した射（第 I 章、3.2.2.2 と 3.2.2.4）であることが直ちに分かる。
したがって、図式 (4.1.7.7) が定める $K(\mathcal A)$ の三角は、定義により
（第 I 章、3.3.1）完全三角である。次に b) を示す。(4.1.6.5) で導入した
記法を用いると、図式

\VerdierSemanticAnchor{ch4d-4-1-7-8}
\begin{equation}
{}^{1}\!Y(q,r)[-1]
\xrightarrow{-{}^{1}\!d(-\infty,q,r)}
{}^{1}\!Y(q)\longrightarrow{}^{1}\!Y(r)\longrightarrow{}^{1}\!Y(q,r),
\tag{4.1.7.8}
\end{equation}

において、複体 ${}^{1}\!Y(r)$ が射 $-{}^{1}\!d(-\infty,q,r)$ の写像錐
として現れることを今示した（図式 (4.1.7.7) で $p=-\infty$ と置く）。
さらに (4.1.1) では、対象 $\mathcal Y(q,r)$ は射

\VerdierSemanticAnchor{ch4d-4-1-7-9}
\begin{equation}
{}^{1}\!Y(q)\longrightarrow{}^{1}\!Y(r).
\tag{4.1.7.9}
\end{equation}

の写像錐として定義された。この状況では（第 I 章、3.2.5.2）、ホモトピー
同値（したがって $K(\mathcal A)$ における同型）

\VerdierSemanticAnchor{ch4d-4-1-7-10}
\begin{equation}
J(q,r)\colon\mathcal Y(q,r)\longrightarrow{}^{1}\!Y(q,r).
\tag{4.1.7.10}
\end{equation}

\VerdierSemanticAnchor{ch4d-4-1-7-proof-end}
が定義されている。このホモトピー同値が
$\mathcal{F}\ell_1(\widetilde{\mathbb Z})$ の指標 $(q,r)$ と二重複体 $X$
に関して関手的であり、スペクトル対象
$\bigl(\mathcal Y(p,q),\delta(p,q,r)\bigr)$ からスペクトル対象
$\bigl({}^{1}\!Y(p,q),{}^{1}\!d(p,q,r)\bigr)$ への標準同型を定めることの
検証は読者に委ねる。
\endgroup
